\begin{abstract}
Hallucinations in large language models (LLMs) pose significant reliability concerns, especially in long-form generations where supported and unsupported claims can appear within the same response. Recent work has shown that LLM internal representations contain useful signals for detecting hallucinations. However, mechanistic
detectors often miss signals distributed across tokens and layers.
Detecting hallucinations in long-form LLM outputs is challing because they often mix supported and unsupported claims, requiring detectors to localize unreliable spans rather than predicting a label for the entire response. Internal representations contain useful hallucination signals, but current mechanistic detectors rely on predetermined tokens or layers. Probes that use fixed layers or tokens miss signals from other positions. Probes that pool activations from the full-response learn global patterns not local ones. Claim extractors and classifiers miss signals before the hallucinated claim. We introduce Temporal Perception Encoder (TPE) that predicts hallucinations based on local signals within a window of tokens and signals from previous tokens.
We achieve this localization by training the model using multi-instance learning over token-centered windows of hidden-state activations. It uses claim or sequence level labes to learn to localize at the token level. Expriments on long-form outputs from biographical generation, open-ended factual QA, and retrieval-augmented QA show that TPE outperforms state-of-the-arts methods.
\end{abstract}
